\subsection*{一步一步建立直观}

\paragraph{1. 先看有向底面，再引入法向量}

两根空间向量张成一个平行四边形。叉积把底面的两个信息合并成一根箭头：模长记录面积，方向垂直于底面并编码输入次序。旋转视角应当只改变屏幕外观，不应改变下面计算得到的面积。



\paragraph{2. 倾斜底面时法线也要跟着转}

把第一根向量抬离平面，红色法向量随底面转动。检查它与两条边的点积都为零，就把“看起来垂直”变成了坐标验证。交换两条边的次序会反转法线，但不会改变平行四边形的普通面积。



\paragraph{3. 共线情形检验定义是否合理}

当两根向量平行时，底面没有面积，叉积应为零。零向量没有唯一方向，因此此时不应继续讨论它指向哪个法线方向。叉积的三维构造把面积与定向联系起来，却不能不加条件地照搬到任意维数。



\subsection*{定义、命题与证明}

\paragraph{定义：叉积}

对三维向量 $u,v$，定义 $u\times v=(u_2v_3-u_3v_2,\ u_3v_1-u_1v_3,\ u_1v_2-u_2v_1)$，方向采用右手定向。

\paragraph{命题：叉积与两个输入都垂直}

对任意 $u,v\in\mathbb R^3$，有 $u\cdot(u\times v)=0$ 且 $v\cdot(u\times v)=0$。叉积可能是零向量。

\begin{proof}

展开第一个点积，得到 $u_1u_2v_3-u_1u_3v_2+u_2u_3v_1-u_2u_1v_3+u_3u_1v_2-u_3u_2v_1$。各项两两抵消，因此和为零。

对第二个点积做同样展开即可。零向量也满足两个点积为零，所以该结论没有宣称平行输入仍产生一个有确定方向的非零法向量。

\end{proof}
